\subsection{p-adic Lie groups.}

Such groups are met for the first time in 1907 in the work of Hensel {[}157 e{]} on the p-adic analytic functions (defined by expansions in entire series). He studies notably the exponential and the logarithm; in spite of the a priori surprising behaviour of the series that define them (for example the exponential series does not converge everywhere), their functional properties remain valid, which supplies a local isomorphism between the additive group and the multiplicative group of \(Q_{p}\) (or, more generally, of every complete ultrametric field of characteristic zero).

It is equally about commutative groups (but not linear this time) that the work of A. Weil {[}330 a{]} and E. Lutz {[}210{]} on the elliptic p-adic curves (1936) is concerned. Other than the arithmetic applications, there can be found the construction of a local isomorphism of the group with the additive group, based on the integration of an invariant differential form. This method can be equally applied to abelian varieties, as is remarked shortly afterwards by C. Chabauty who uses it without any more explanation in order to prove a particular case of the ``Mordell conjecture'' {[}61{]}.

From this moment, it was clear that the local theory of Lie groups can be applied almost without change to the p-adic case. The fundamental theorems of the ``dictionary'' Lie groups-Lie algebras are established in 1942 in the thesis of R. Hooke {[}166{]}, a pupil of Chevalley; this work contains also the p-adic analogue of the theorem of E. Cartan on the closed subgroups of real Lie groups.

More recently, M. Lazard {[}195 b{]} develops a more precise form of the ``dictionary'' for compact analytic groups over \(\mathbf{Q}_p\) . He shows that the existence of an analytic \(p\) -adic structure on a compact group \(G\) is tightly linked to that of certain filtrations on \(G\) , and gives various applications of it (for example to the cohomology of \(G\) ). One of Lazard's tools is an improvement of Dynkin's results on the convergence of the \(p\) -adic Hausdorff series {[}195 a{]}.

\section{FREE LIE ALGEBRAS.}

It remains for us to speak of a series of works on Lie algebras where the link with the theory of Lie groups is very tenuous; these researches have on the other hand important applications in the theory of ``abstract'' groups and more especially nilpotent groups.

The origin of this is the work of P. Hall {[}144{]}, which appeared in 1932. There is however there no question of Lie algebras: P. Hall has in view the study of a certain class of p-groups, those that he calls ``regular''. But that leads him to examine in detail iterated commutators and the lower central series of a group; he establishes on this occasion a variant of the Jacobi identity, as well as the ``Hall formula''

\[
(x y) ^ {n} = x ^ {n} y ^ {n} (x, y) ^ {n (n - 1) / 2} \dots
\]

Soon after (in 1935-1937) the fundamental works of W. Magnus {[}215 a and b{]} and E. Witt {[}337 b{]} appear. In {[}215 a{]} Magnus uses the same algebra of formal series \(\hat{A}\) as Hausdorff (since called the ``Magnus algebra''); he embeds the free group \(F\) in it and uses the natural filtration of \(\hat{A}\) in order to obtain a decreasing sequence \((F_n)\) of subgroups of \(F\) ; it is one of the first examples of filtration. He conjectures that the \(F_n\) coincide with the terms of the lower central series of \(F\) . This conjecture is proved in his second memoir {[}215 b{]}; it is also there that he explicitly makes the link between his ideas and those of P. Hall, and that he defines the free Lie algebra \(L\) (as a sub algebra of \(A\) ) of which he shows in substance that it is identified with the graded \(F\) . In {[}337 b{]}, Witt completes this result on several points. He shows notably that the enveloping algebra of \(L\) is a free associative algebra and deduces from that immediately the rank of the homogeneous components of \(L\) (``Witt formulae'').

As for the determination of the basis of L known by the name of ``Hall basis'', it seems that it only appears for the first time in 1950, in a note of M. Hall {[}143{]}, although it is implicit in the works of P. Hall and W. Magnus quoted above.

\chapter{GROUPS GENERATED BY REFLECTIONS; ROOT SYSTEMS.}

The groups considered in this Note appeared in connection with various questions of Geometry, Analysis and the Theory of Lie groups, sometimes in the form of permutation groups, sometimes in the from of groups of displacements in Euclidean or hyperbolic geometry, and these various points of view have only been co-ordinated in recent times.

Historically, the beginnings of the theory are a good deal older than the introduction of the concept of a group: it takes its source indeed in the studies on the ``regularity'' or the ``symmetries'' of geometrical figures, and notably in the determination of regular polygons and polyhedra (going back no doubt to the Pythagoreans), which constitute the crowning achievement of the Elements of Euclid, and one of the most admirable creations of the Greek genius. Later, notably with the Arab authors of the high Middle Ages, then with Kepler, the beginnings of a mathematical theory of regular ``tilings'' of the plane or the sphere by polygons congruent in pairs (but not necessarily regular) appear, no doubt linked to the origin of the different types of ornament thought up by the ancient and Arab civilisations (that can be considered with good cause as an authentic part of the mathematics developed by these civilisations {[}291{]}).

Towards 1830-1840, studies in crystallography (Hessel, Bravais, Möbius) lead to considering a problem that is exactly that of the determination of the finite groups of displacements in Euclidean space of 3 dimensions, although the authors quoted above do not use yet the language of the theory of groups; this latter hardly comes into use until around 1860, and it is in the context of the classification of groups that Jordan, in 1869 {[}174 b{]}, determines the discrete subgroups of displacements in \(\mathbf{R}^3\) which preserve orientation (and more generally, all the closed subgroups of the groups of displacements preserving orientation).

Until the last years of the XIXth century, this stream of ideas is developed in many directions, of which the most marked are the following:

\begin{enumerate}
\def\labelenumi{\arabic{enumi}.}
\tightlist
\item
  Conforming to a tendency that appears very early on in the theory of finite groups, ``presentations'' of finite groups of displacements by generators and relations of a simple type are sought. It is thus that Hamilton, already in 1856 {[}145 b{]}, proves that the finite groups of rotations in Euclidean space
\end{enumerate}

\(R^{3}\) are generated by two generators S, T linked by the relations \(S^{p} = T^{q} = (ST)^{3} = 1\) for appropriate values of p and q.

\begin{enumerate}
\def\labelenumi{\arabic{enumi}.}
\setcounter{enumi}{1}
\item
  Discrete groups of displacements may or may not contain reflections. Already in 1852, Möbius determines substantially the finite groups of displacements in spherical geometry generated by reflections (which is equivalent to the same problem for finite groups of Euclidean displacements in \(\mathbf{R}^3\) ); he finds that, with the exception of the cyclic groups, such a group has as fundamental domain a spherical triangle having angles of the form \(\pi / p, \pi / q, \pi / r\) , where \(p, q, r\) are three integers \(>1\) such that \(\frac{1}{p} + \frac{1}{q} + \frac{1}{r} > 1\) ({[}223{]}, v. II, pp.~349-360 and pp.~561-708). He also notices that these groups contain all the finite groups of displacements as subgroups.
\item
  This last stream of ideas finds a new amplification when following work of Riemann and Schwarz on hypergeometric functions and conformal representations, the study of ``tilings'' of the complex plane or of the half-plane by means of figures bounded by arcs of a circle is started; Klein and Poincaré make it the foundation of the theory of ``automorphic functions'' and recognise in it (for the case of arcs of a circle orthogonal to a fixed straight line) a problem equivalent to that of the search for the discrete subgroups of the group of displacements of the non-Euclidean hyperbolic plane (identified with the ``Poincaré half-plane'') {[}118{]}.
\item
  The notions of regular polyhedron and of the tiling of \(R^{3}\) by such polyhedra are extended to all the Euclidean spaces \(R^{n}\) by Schläfli, in work that goes back to the neighbourhood of 1850, but was only published much later and remained ignored for a long time ({[}273{]}, v. I, pp.~167-387); he determines completely the regular ``polytopes'' in each \(R^{n}\) , the group of displacements leaving invariant such a polytope, and a fundamental domain of this group, which, as in the case n = 3 studied by Möbius, is a ``room'' whose trace on the sphere \(S_{n-1}\) is a spherical simplex. However, he does not tackle the inverse problem of the search for the finite displacement groups generated by reflections in \(R^{n}\) ; this problem will only be solved much later, by Goursat {[}132{]} for n = 4, and, for arbitrary n, its solution will have to wait for the work of E. Cartan ({[}52 a{]}, v. I₂, pp.~1003-1020) and of Coxeter {[}70 c{]}, to which we will return later on.
\end{enumerate}

Around 1890, with the first work of Killing and of E. Cartan on Lie groups, begins a new stream of ideas that for a long time will develop without links to the preceding work. Killing {[}180{]} and Cartan ({[}52 a{]}, v. I \(_{1}\) , pp.~137-287), in their study of the structure of complex semisimple Lie algebras, will plan straightaway a primordial role for certain linear forms \(\omega_{\alpha}\) on a ``Cartan sub algebra'' \(\mathfrak{h}\) of such a Lie algebra \(\mathfrak{g}\) ; they are the ``roots'' relative to \(\mathfrak{h}\) , so called because with Killing they appear as the roots of the characteristic equation \(\det(\mathrm{ad}_{\mathfrak{g}}(x) - T) = 0\) , considered as functions of \(x \in \mathfrak{h}\) . The properties of these ``roots'' established by Killing and Cartan reduce to affirming that, in the actual geometric language, they form a ``reduced system of roots'';

they show then that the classification of the complex semisimple Lie algebras is reduced to that of associated ``systems of roots'', which itself reduces to the determination of certain matrices with integer coefficients (later called ``Cartan matrices''). Killing and Cartan also bring out, for every root \(\omega_{\alpha}\) , the existence of an involutive permutation \(S_{\alpha}\) of the set of roots; they use in an essential way the transformation \(C = S_{\alpha_{1}}S_{\alpha_{2}} \ldots S_{\alpha_{l}}\) , the product of permutations associated with l roots making up a fundamental system (a transformation today called a ``Coxeter transformation''); they even extend this permutation into a linear transformation of the vector space generated by the fundamental roots \(\omega_{\alpha_{i}}(1 \leq i \leq l)\) , and study its eigenvalues ({[}180{]}, II, p.~20; {[}52 a{]}, v. I \(_{1}\) , p.~58). But neither Killing, nor at first Cartan, seem to have thought of considering the group \(G'\) generated by the \(S_{\alpha}\) ; and when Cartan, a little later ({[}52 a{]}, v. I \(_{1}\) , pp.~293-353), determines the Galois group G of the characteristic equation

\[
\det (\mathrm{ad} _ {\mathfrak {g}} (x) - T) = 0
\]

of a ``general element'' \(x \in h\) , he studies it at first without bringing in the \(S_{\alpha}\) ; thirty years later, already under the influence of the work of H. Weyl, he proves ({[}52 a{]}, v. I \(_{1}\) , pp.~555-568) that G has as a normal subgroup the group \(G'\) and determines in all cases the structure of the quotient group \(G/G'\) which (for a simple algebra g) is or order 1 or 2 except for the type \(D_{4}\) where it is isomorphic to \(S_{3}\) ; it is also on this occasion that he interprets \(G'\) as the group induced by the inner automorphisms of a complex semisimple Lie algebra, leaving fixed a Cartan subalgebra.

The work of H. Weyl, to which we have just alluded, is that which inaugurated the geometric interpretation of the group \(G'\) (since called the ``Weyl group'' of g); in the same way as Killing and Cartan had done it for the transformation C, he had the idea of considering the \(S_{\alpha}\) as reflections in the vector space of linear forms on h. It is also in the memoir of H. Weyl {[}331 b{]} that the fundamental domain of the ``affine Weyl group'' is seen to appear (without besides the link with the ``Weyl group'' \(G'\) being clearly indicated); Weyl uses it in order to prove that the fundamental group of a compact semisimple group is finite, a crucial point in his proof of the complete reducibility of the linear representations of a complex semisimple Lie algebra. Soon after, E. Cartan brings to fruition the synthesis of the global points of view of H. Weyl, of his own theory of real or complex semisimple Lie algebras, and of the theory of symmetric Riemann spaces that he was building at this time. In the memoir ({[}52 a{]}, v. I₂, pp.~793-840), he completes the determination of the fundamental polytopes of the Weyl group and of the affine Weyl group, and introduces the weight systems and radical weight systems; he extends this discussion to symmetric spaces, and meets thus notably the first examples of non reduced root systems ({[}52 a{]}, v. I₂, pp.~1003-1010). Finally the article ({[}52 a{]}, v. I₂, pp.~867-989) gives the first proof of the fact that every finite group generated by reflections in \(R^{n}\) and irreducible has a fundamental domain having as trace on \(S_{n-1}\) a spherical simplex; it is also in this work that he proves the uniqueness of the largest root (under an arbitrary lexicographic ordering on the system of roots) by geometrical considerations.

A little later, van der Waerden {[}317 b{]}, relying on the memoir of H. Weyl, shows that the classification of the complex semisimple Lie algebras is equivalent to that of the reduced root systems, which he carries out by elementary geometric considerations (whereas, with Killing and Cartan, this classification is a result of complicated calculations with determinants). At about the same time, Coxeter determines explicitly all the finite irreducible groups of Euclidean displacements which are generated by reflections {[}70 c{]}; he completes thus the results of the memoir ({[}52 a{]}, v. I₂, pp.~793-840) of E. Cartan, who had only determined the ``crystallographic'' groups (i. e. associated to a root system, or also susceptible of being embedded in an infinite discrete group of displacements). The following year {[}70 d{]}, Coxeter shows that the finite groups generated by reflections are the only finite groups (up to isomorphism) admitting a presentation by involutive generators \(R_i\) subject to relations of the form \((R_i R_j)^{m_{ij}} = 1\) ( \(m_{ij}\) integers), whence the name of ``Coxeter groups'' given since to groups (finite or not) admitting such a presentation.

The first link between the two streams of research that we have described above seems to have been established by Coxeter {[}70 e{]}, then by Witt {[}337 c{]}. They notice that the infinite irreducible groups of Euclidean displacements generated by reflections correspond bijectively (up to isomorphism) to complex simple Lie algebras. Witt gives a new determination of discrete groups of this type, and extends also the theorem of Coxeter {[}70 d{]} recalled above by characterising as well the Coxeter groups isomorphic to infinite discrete groups of Euclidean displacements. This result, and the fact that the analogous groups in hyperbolic geometry are also Coxeter groups, \(^{1}\) have led to the open tackling of the study of these latter, first by putting the emphasis on a geometric realisation ({[}71{]}, {[}308 a{]}), then, following J. Tits {[}308 b and c{]} in a purely algebraic framework.

Starting with the work of Witt, the theory of semisimple Lie groups and that of discrete groups generated by reflections will not cease from reacting extremely fruitfully the one on the other. Already from 1941, Stiefel {[}296{]} remarks that the Weyl groups are exactly the finite groups generated by reflections, and that leave a network invariant. Chevalley {[}62 e{]} and Harish-Chandra {[}149 a{]} give in 1948-51 proofs a priori of the bijective correspondence between ``crystallographic'' groups and complex semisimple Lie algebras; it was not possible until then to verify separately this correspondence for each type of simple Lie algebra.

Around 1950, it is noticed on the other hand that the polynomials invariant under the Weyl group play an important role in the theory of linear representations of infinite dimension {[}149 a{]} and in the topology of Lie groups.

On his side, Coxeter {[}70 g{]}, taking up again the study of the transformation C, product of the fundamental reflections of a finite group W generated by reflections, notices (by means of a separate examination of each type) that the algebra of polynomials invariant under W is generated by algebraically independent elements, whose degrees are linked in a simple way to the eigenvalues of C. Proofs a priori of these results were subsequently given by Chevalley {[}62 f{]} for the first, and by Coleman {[}68{]} and Steinberg {[}292{]} for the second.

With the work of A. Borel on linear algebraic groups {[}30{]} new developments in the theory of Lie groups start which were to lead to a notable enlargement of this latter. A. Borel brings out the importance of the maximal connected soluble subgroups (since called ``Borel subgroups'') in a Lie group, and makes of them the principal tool for transposing a large part of the classical theory to algebraic groups over an algebraically closed field (without however obtaining yet a classification of the simple algebraic groups \(^{2}\) ). The Borel subgroups (in the case of real or complex classical groups) had already turned up some years previously in the work of Gelfand and Neumark on the representations of infinite dimension; and in 1954, F. Bruhat had discovered the remarkable fact that, for the classical simple groups, the decomposition of the group into double cosets over a Borel group is indexed in a canonical way by the Weyl group {[}40{]}. This result was subsequently extended to all real and complex semisimple groups by Harish-Chandra {[}149 b{]}. On the other hand, in 1955, Chevalley {[}62 g{]} had succeeded in associating with every complex semisimple Lie algebra g and with every commutative field k, a group of matrices with coefficients in k, possessing a Bruhat decomposition; and he used this last fact to show that, with a small number of exceptions, the group thus defined was simple (in the sense of the theory of abstract groups). He ``explained'' thus the coincidence, already noticed since Jordan and Lie, between the simple Lie groups (in the sense of the theory of Lie groups) of types A, B, C, D and the classical simple groups defined in a purely algebraic way over an arbitrary field (a coincidence that until then could not be extended except to the exceptional types \(G_{2}\) and \(E_{6}\) by Dickson {[}88c and d{]}). In particular, by taking a finite field k, the construction of Chevalley supplied, for each type of complex simple Lie algebra, a family of finite simple groups, containing a large part of the then known finite simple groups, as well as three new series (corresponding to the types of simple Lie algebras \(F_{4}, E_{7}\) and \(E_{8}\) ). Soon after, by various procedures, using modifications of the methods of Chevalley, several authors (Hertzig, Suzuki, Ree, Steinberg and Tits) on the one hand showed that one can obtain in an analogous way the other finite simple groups known at that time, with the exception of the alternating groups and the Mathieu groups, and on the other hand constructed other series of new finite simple groups (cf.~{[}54{]}).

At almost the same time, Chevalley {[}62 h{]}, still using the technique of Bruhat decompositions, joined to a key result on the normaliser of a Borel subgroup, took up again the study of linear algebraic groups and reached the result that over an algebraically closed field k of arbitrary characteristic, the theory of semisimple linear algebraic groups \(^{3}\) leads essentially to the same types as in the Killing-Cartan classification for k = C. Following on, J. Tits {[}308 a and b{]}, analysing the methods of Chevalley, reached an axiomatised version (the ``BN-pairs'') of the Bruhat decompositions, in a remarkably flexible form, only bringing in the structure of the group; it is this notion that is at present known under the name ``Tits system''. All the simple groups (in the various meanings of the word) of which we have been speaking above are canonically supplied with Tits systems, and Tits himself {[}308 c{]} has proved that the existence of such a system in an abstract group G, linked with a few supplementary properties from pure group theory, allows the proof that G is simple, a theorem which covers the majority of the proofs of simplicity given until then for these groups. In collaboration with A. Borel, he has on the other hand generalised the results of Chevalley of {[}62 h{]}, by showing the existence of Tits systems in the group of rational points of a semisimple linear algebraic group over an arbitrary field {[}31{]}.

All the Tits systems occurring in these questions have a finite Weyl group. Another category of examples was discovered by Iwahori and Matsumoto {[}170{]}; they have shown that if, in Chevalley's construction of {[}62 g{]}, k is a p-adic field, then the group obtained has a Tits system whose Weyl group is the affine Weyl group of the complex semisimple Lie algebra with which one started. This result has just been extended by Bruhat and Tits {[}41{]} to all semisimple algebraic groups over a local field.

\specialchapter{Bibliography}

1 N. H. Abel. Oeuvres, 2 vol.~ed.~Sylow and Lie, Christiana, 1881.

2a I. Ado. Note on the representation of continuous and finite groups by means of linear substitutions (in Russian), Bull. Phys. Math. Soc. Kazan, v. VII (1935), pp.~3-43.

2b I. Ado. The representation of Lie algebras by matrices (in Russian), Uspehi Mat. Nauk., v. II (1947), pp.~159-173 (English trans.: Amer. Math. Soc. Transl., (1), vol.~9, pp.~308-327).

3 A. D. Alexandroff. Additive set functions in abstract spaces, Mat. Sbornik: I (chap.~1), v. VIII (1940), pp.~307-348; II (chap.~2 and 3), v. IX (1941), pp.~563-628; III (chap 4 to 6), v. XIII (1943), pp.~169-238.

4 P. Alexandroff-H. Hopf. Topologie I, Berlin (Springer), 1935.

5a Archimedis Opera quae quidem exstant omnia, nunc primus et gr. et lat. edita \ldots. Basilae, Jo. Hervagius, 1 vol.~in-fol., 1544.

5b Archimedis Opera Omnia. 3 vols., ed.~J. L. Heiberg, 2nd ed., Leipzig (Teubner), 1913-1915.

5c Les Oeuvres complètes d'Archimède. Transl. P. Ver Ecke, Paris-Brussels (Desclée-de Brouwer), 1921.

6 The works of Aristotle. Translated under the editorship of W. D. Ross, Oxford, 1928 sqq.

7a E. Artin. Galois theory \ldots, Notre Dame, 1946.

7b E. Artin. Ueber die Zerlegung definiter Funktionen in Quadraten, Abh. Math. Sem. Univ. Hamburg, v. V (1927), pp.~100-115.

7c E. Artin. Zur Theorie der hypercomplexen Zahlen, Abh. Math. Sem. Univ. Hamburg, v. V (1927), pp.~251-260.

7d E. Artin. Einführung in die Theorie der Gammafunktion, Leipzig (Teubner), 1931.

8a E. Artin und O. Schreier. Algebraische Konstruktion reeler Körper, Abh. Math. Sem. Univ. Hamburg, v. V (1927), pp.~85-99.

8b E. Artin und O. Schreier. Eine Kennzeichnung der reell abgeschlossenen Körper, Abh. Math. Sem. Univ. Hamburg, v. V (1927), pp.~225-231.

9a C. Arzelà. Sulla integrabilità di una serie di funzioni, Rendic. Acc. dei Lincei (4), v. I (1885), pp.~321-326.

9b C. Arzelà. Sulla integrazione per serie. Rendic. Acc. dei Lincei (4), v. I (1885), pp.~532-537 and 566-569.

10 G. Ascoli. Le curve limiti di una varietà data di curve, Mem. Acc. dei Lincei (3), v. XVIII (1883), pp.~521-586.

11a R. Baire. Leçons sur les fonctions discontinues, Paris (Gauthier-Villars), 1905.

11b R. Baire. Sur les fonctions de variables réelles, Ann. di Mat. (3), v. III (1899), pp.~1-123.

12 R. Baire, E. Borel, J. Hadamard, H. Lebesgue. Cinq lettres sur la théorie des ensembles, Bull. Soc. Math. de France, v. XXXIII (1905), pp.~261-273.

13 H. F. Baker. Alternants and continuous groups, Proc. Lond. Math. Soc. (2), v. III (1905), pp.~24-47.

14a S. Banach. Théorie des opérations linéaires, Warszawa, 1932.

14b S. Banach. Sur le problème de la mesure, Fund. Math., v. IV (1923), pp.~7-33.

14c S. Banach. Sur les fonctionnelles linéaires, Studia Math., v. I (1929), pp.~211-216 and 223-239.

15 S. Banach and H. Steinhaus. Sur le principe de condensation des singularités, Fund. Math., v. IX (1927), pp.~50-61.

16a I. Barrow. Lectiones Geometricae \ldots, Londini, 1670.

16b I. Barrow. Mathematical Works, ed.~W. Whewell, Cambridge, 1860.

17a O. Becker. Eudoxos-Studien, Quellen und Studien zur Geschichte der Math., Abt. B: Studien, v. II (1933), pp.~311-333 and 369-387, and v. III (1936), pp.~236-244 and 370-388.

17b O. Becker. Die Lehre vom Geraden und Ungeraden im Neunten Buch der Euklidischen Elementen, Quellen und Studien zur Geschichte der Math., Abt. B : Studien, v. III (1936), pp.~533-553.

18a E. Beltrami. Saggio di interpretazione della geometrica non-euclidean, Giorn. di Mat., v. VI (1868), pp.~284-312.

18b E. Beltrami. Teoria fondamentale degli spazii di curvatura costante, Ann. di Mat. (2), v. II (1868-69), pp.~232-255.

19a Jakob Bernoulli. Opera, 2 vol;., Geneva (Cramer-Philibert), 1744.

19b Jakob Bernoulli. Ars Conjectandi, Basel, 1713 (transl. S. Vastel, Paris, 1801).

20a Johann Bernoulli. Opera Omnia, 4 vol., Lausanne-Geneva (m. M. Bousquet), 1742.

20b Johann Bernoulli. Der erste Integralrechnung (Ostwald's Klassiker, no 194), Leipzig-Berlin (Engelmann), 1914.

20c Johann Bernoulli. Die Differentialrechnung (Ostwald's Klassiker, no 211), Leipzig (Akad. Verlag), 1924.

21 E. Bezout. Théorie générale des équations algébriques, Paris, 1779.

22a G. Birkhoff. Continuous groups and linear spaces, Rec. Math. Moscow, v. I, (1936), pp.~635-642.

22b G. Birkhoff. Representability of Lie algebras and Lie groups by matrices, Ann. of Math., v. XXXVIII (1937), pp.~526-532.

23 J. M. Bochenski. Ancient formal logic, Studies in Logic, Amsterdam (North-Holland Publ. Co.), 1951.

24a S. Bochner. Integration von Funktionen deren Werte die Elemente eines Vektorraumes sind, Fund. Maht., v. XX (1933), 262-276.

24b S. Bochner. Harmonic Analysis and the Theory of Probability, Berkeley (Univ. of California Press), 1960.

25 P. Böhner. Medieval logic, an outline of its development from 1250 to ca. 1400, Chicago, 1952.

26 H. Bohr und J. Mollerup. Laerebog i matematisk Analyse, v. III, Kopenhagen, 1922.

27a B. Bolzano. Oeuvres, 5 vol., Prague, 1930-1948.

27b B. Bolzano. Paradoxien des Unendlichen, Leipzig, 1851 (English trans., New Haven (Yale Univ. Press), 1950).

27c B. Bolzano. Rein Analytischer Beweis der Lehrsatzes, dass zwischen zwei Werthen, die ein entgegengesetzes Resultat gewahren, wenigstens eine reelle Wurzel liegt (Ostwald's Klassiker, no 153), Leipzig, 1905.

28a R. Bombelli. L'Algebra, Bologna (G. Rossi), 1579.

28b R. Bombelli, L'Algebra, libri IV e V, ed E. Bortolotti, Bologna (Zanichelli), 1929.

29 G. Boole. Collected logical works, 2 vol., ed.~P. Jourdain, Chicago- London, 1916.

30 A. Borel. Groupes linéaires algébriques, Ann. of Math., v. LXIV (1956), pp.~20-80.

31 A. Borel and J. Tits. Groupes réductifs, Publ. math. I. H. E. S., no 27 (1965), pp.~55-150.

32a E. Borel. Leçons sur la théorie des fonctions, Paris (Gauthier-Villars), 1898.

32b E. Borel. Les probabilités dénombrables et leurs applications arithmétiques, Rend. Circ. Mat. Palermo, v. XXVII (1909), pp.~286-310.

33 H. Bosmans. Sur le ``libro del Algebra'' de Pedro Nuñez, Bibl. Math. (3), v. VIII (1907-08), pp.~154-169.

34a R. Brauer. Über Systeme hypercomplexer Zahlen, Math. Zeitschr., v. XXX (1929), pp.~79-107.

34b R. Brauer. Eine Bedingung für vollständige Reduzibilität von Darstellungen gewöhnlicher und infinitesimaler Gruppen, Math. Zeitschr., v. XLI (1936), pp.~330-339.

35 A. Bravais. Mémoires sur les polyèdres de forme symétrique, Journ. de Math. (1), v. XIV (1849), pp.~141-180.

36 H. Briggs. Arithmetica logarithmica, London, 1624.

37 A. Brill und M. Noether. Ueber algebraische Funktionen, Math. Ann., v. VII (1874), pp.~269-310.

38 Lord Brouncker. The Squaring of the Hyperbola by an infinite series of Rational Numbers \ldots, Phil. Trans., v. III (1668), pp.~645-649.

39a L. E. J. Brouwer. Intuitionism and formalism, Bull. Amer. Math. Soc., v. XX (1913), pp.~81-96.

39b L. E. J. Brouwer. Zur Begründung des intuitionistischen Mathematik, Math. Ann., v. XCIII (1925), pp.~244-257, v. XCV (1926), pp.~453-473, v. XCVI (1926), pp.~451-458.

40 F. Bruhat. Représentations induites des groupes de Lie semi-simples complexes, C. R. Acad. Sci., v. CCXXXVIII (1954), pp.~437-439.

41 F. Bruhat and J. Tits. Groupes algébriques simples sur un corps local, Proc. Conf. on Local Fields, pp.~23-36, Berlin (Springer), 1967.

42 C. Burali-Forti. Sopra un theorema del Sig. G. Cantor, Atti Acad. Torino, v. XXXII (1896-97), pp.~229-237.

43 H. Burkhardt. Die Anfänge der Gruppentheorie und Paolo Ruffini, Zeitschr. für Math. und Phys., v. XXXVII (1892), Suppl., pp.~121-159.

44a W. Burnside. On the condition of reducibility for any group of linear substitutions, Proc. Lond. Math. Soc., v. III (1905), pp.~430-434.

44b W. Burnside. Theory of groups of finite order, 2nd ed., Cambridge (University Press), 1911.

45 W. H. Bussey. The origin of mathematical induction. Amer. Math. Monthly, v. XXIV (1917), pp.~199-207.

46a J. E. Campbell. On a law of combination of operators bearing on the theory of continuous transformation groups, Proc. Lond. Math. Soc., (1), v. XXVIII (1897), pp.~381-390.

46b J. E. Campbell. On a law of combination of operators (second paper), Proc. Lond. Math. Soc., (1), v. XXIX (1898), pp.~14-32.

47 G. Cantor. Gesammelte Abhandlungen, Berlin (Springer), 1932.

48 G. Cantor, R. Dedekind. Briefwechsel, ed.~J. Cavaillès-E. Noether, Actual. Scient. et Ind., no 518, Paris (Hermann), 1937.

49 C. Caratheodory. Vorlesungen über reelle Funktionen, Leipzig-Berlin (Teubner), 1918.

50 H. Cardano. Opera, 10 vol., Lyon, 1663.

51 L. Carnot. Géométrie de Position, Paris, 1803.

52a E. Cartan. Oeuvres complètes, 6 vol.~(3 parts), Paris (Gauthier-Villars), 1953-55.

52b E. Cartan. Leçons sur la théorie des spineurs (Actual. Scient. et Ind., nos 643 and 701), Paris (Hermann), 1938.

53 H. Cartan. Théorie des filtres, C. R. Acad. Sci., v. CCV (1937), pp.~595-598; Filtres et ultrafiltres, ibid., pp.~777-779.

54 R. Carter. Simple groups and simple Lie algebras, Journ. Lond. Math. Soc., v. XL (1965), pp.~193-240.

55 H. Casimir und B. L. Van der Waerden. Algebraischer Beweis der vollständigen Reduzibilität der Darstellungen halbeinfacher Liescher Gruppen, Math. Ann., v. CXI (1935), pp.~1-12.

56a A.-L. Cauchy. Oeuvres complètes, 26 vol., (2 series), Paris (Gauthier-Villars), 1882-1958.

56b Leçons de calcul différentiel et de calcul intégral, rédigées principalement d'après les méthodes de M. A.-L. Cauchy, by l'Abbé Moigno, v. II, Paris, 1844.

57a B. Cavalieri. Geometria indivisibilibus continuorum quadam ratione promota, Bononiae, 1635 (2nd ed., 1653).

57b B. Cavalieri. Exercitationes geometricae sex, Bononiae, 1647.

58 A. Cayley. Collected Mathematical Papers, 13 vol., Cambridge (University Press), 1889-1898.

59 L. Cesari. Surface area, Princeton, 1954.

60a M. Chasles. Note sur les propriétés générales du système de deux corps, Bull. de Férussac, v. XIV (1830), pp.~321-326.

60b M. Chasles. Aperçu historique sur l'origine et le développement des méthodes en géometrie, Brussels, 1837.

61 C. Chabauty. Sur les points rationnels des courbes algébriques de genre supérieur à l'unité, C. R. Acad. Sci., v. CCXII (1941), pp.~882-884.

62a C. Chevalley. The algebraic theory of spinors, New-York (Columbia University Press), 1954.

62b C. Chevalley. Généralisation de la théorie du corps de classes pour les extensions infinies, Journ. de Math., (9), v. XV (1936), pp.~359-371.

62c C. Chevalley. On the theory of local rings, Ann. of Math., v. XLIV (1943), pp.~690-708.

62d C. Chevalley. Theory of Lie groups, Princeton (University Press), 1946.

62e C. Chevalley. Sur la classification des algèbres de Lie simples et de leurs représentations, C. R. Acad. Sci., v. CCXXVII (1948), pp.~1136-1138.

62f C. Chevalley. Invariants of finite groups generated by reflections, Amer. Journ. of Math., v. LXXVII (1955), pp.~778-782.

62g C. Chevalley. Sur certains groupes simples, Tohōku Math. Journ., (2), v. VII (1955), pp.~14-66.

62h C. Chevalley. Classification des groupes de Lie algébriques, 2 vol., Paris (Inst. H. Poincaré), 1956-58.

63 G. Choquet. Theory of capacities, Ann. Inst. Fourier, v. V (1953-54), pp.~131-295.

64 N. Chuquet. Le Triparty en la Science des Nombres, ed.~A. Marre, Bull. bibl. storia math., v. XIII (1880), pp.~555-659 and 693-814.

65 W. K. Clifford. Mathematical Papers, London (Macmillan), 1882.

66 I. Cohen and A. Seidenberg. Prime ideals and integral dependence, Bull. Amer. Math. Soc., v. LII (1946), pp.~252-261.

67 P. J. Cohen. The independence of the continuum hypothesis, Proc. Nat. Acad. Sci. U. S. A., v. L (1963), pp.~1143-1148 and v. LI (1964), pp.~105-110.

68 A. J. Coleman. The Betti numbers of the simple groups, Can. Journ. of Math., v. X (1958), pp.~349-356.

69a L. Couturat. La logique de Leibniz d'après des documents inédits, Paris (Alcan), 1901.

69b L. Couturat. La Philosophie des mathématiques de Kant, Revue de Métaph. et de Morale, v. XII (1904), pp.~321-383.

70a H. S. M. Coxeter. Groups whose fundamental regions are simplexes, Journ. Lond. Math. Soc., v. VI (1931), pp.~132-136.

70b H. S. M. Coxeter. The polytopes with regular prismatic figures, Proc. Lond. Math. Soc., (2), v. XXXIV (1932), pp.~126-189.

70c H. S. M. Coxeter. Discrete groups generated by reflections, Ann. of Math., v. XXXV, (1934), pp.~588-621.

70d H. S. M. Coxeter. The complete enumeration of finite groups of the form \(R_{i}^{2} = (R_{i}, R_{j})^{k_{ij}} = 1\) , Journ. Lond. Math. Soc., v. X (1935), pp.~21-25.

70e H. S. M. Coxeter in H. Weyl. The structure and representation of continuous groups (Inst. for Adv. Study, mimeographed notes by N. Jacobson and R. Brauer, 1934-35): Appendix.

70f H. S. M. Coxeter. Regular polytopes, New York (Macmillan), 1948 (2nd ed., 1963).

70g H. S. M. Coxeter. The product of generators of a finite group generated by reflections, Duke Math. Journ., v. XVIII (1951), pp.~765-782.

71 H. S. M. Coxeter and W. O. J. Moser. Generators and relations for discrete groups, Ergebn. der Math., Neue Folge, Bd. 14, Berlin (Springer), 1957 (2nd ed., 1965).

72 G. Cramer. Introduction à l'analyse des lignes courbes, Geneva (Cramer and Philibert), 1750.

73 H. Curry. Outlines of a formalist philosophy of mathematics, Amsterdam (North-Holland Publ. Co.), 1951.

74 M. Curtze. Über die Handschrift ``Algorismes proportionum magistri Nicolay Orem'', Zeitschr. für Math. und Phys., v. XIII, Supplem. (1868), pp.~65-79 and pp.~101-104.

75a J. d'Alembert. Encyclopédie, Paris, 1751-1765.

75b J. d'Alembert. Sur les principes métaphysiques de Calcul infinitésimal, Mélanges de litt., d'hist. et de philosophie, new ed., v. V, Amsterdam (1768), pp.~207-219.

75c J. d'Alembert. Misc. Taur., v. III (1762-65), p.~381.

76a P. J. Daniell. A general form of integral, Ann. of Math. (2), v. XIX (1918), pp.~279-294.

76b P. J. Daniell. Integrals in an infinite number of dimensions. Ann. of Math., (2), v. XX (1919), pp.~281-288.

76c P. J. Daniell. Stieltjes-Volterra products, Congr. Intern. des Math., Strasbourg, 1920, pp.~130-136.

76d P. J. Daniell. Functions of limited variation in an infinite number of dimensions, Ann. of Math., (2), v. XXI (1919-20), pp.~30-38.

77 G. Darboux. Mémoire sur les fonctions discontinues, Ann. Ec. Norm. Sup. (2), v. IV (1875), pp.~57-112.

78 B. Datta and A. N. Singh. History of Hindu Mathematics, 2 vol., Lahore (Motilal Banarsi Das), 1935-38.

79 R. Dedekind. Gesammelte mathematische Werke, 3 vol., Braunschweig (Vieweg), 1932.

80 R. Dedekind und H. Weber. Theorie der algebraischen Funktionen einer Veränderlichen, J. de Crelle, v. XCII (1882), pp.~181-290.

81 M. Dehn. Über den Rauminhalt, Math. Ann., v. LV (1902), pp.~465-478.

82 Ch. de la Vallée Poussin. Intégrales de Lebesgue, Fonctions d'ensembles, Classes de Baire, Paris (Gauthier-Villars), 1916 (2nd ed., 1936).

83a A. de Morgan. On the syllogism (III), Trans. Camb. Philos. Soc., v. X (1858), pp.~173-230.

83b A. de Morgan. On the syllogism (IV) and on the logic of relations, Trans. Camb. Philos. Soc., v. X (1860), pp.~331-358.

84 G. Desargues. Oeuvres, ed.~Pourra, v. I, Paris (Leiber), 1864.

85a R. Descartes. Oeuvres, ed.~Ch. Adam and P. Tannery, 13 vol., Paris (L. Cerf), 1897-1913.

85b R. Descartes. Geometria, latin transl. of Fr.~van Schooten, 2nd ed., 2 vol., Amstardam (Elzevir), 1659-61.

86 J. A. de Séguier. Théorie des groupes finis. Éléments de la théorie des groupes abstraits, Paris (Gauthier-Villars), 1904.

87 M. Deuring. Algebren (Erg. der Math., Bd. 4), Berlin (Springer), 1937.

88a L. E. Dickson. Linear associative algebras and abelian equations, Trans. Amer. Math. Soc., v. XV (1914), pp.~31-46.

88b L. E. Dickson. Theory of linear groups in an arbitrary field, Trans. Amer. Math. Soc., v. II (1901), pp.~363-394.

88c L. E. Dickson. A new system of simple groups, Math. Ann., v. LX (1905), pp.~137-150.

88d L. E. Dickson. A class of groups in an arbitrary realm connected with the configuration of the 27 lines on a cubic surface, Quart. Journ. Pure and App. Math., v. XXXIII (1901), pp.~145-173 and v. XXXIX (1908), pp.~205-209.

89 H. Diels. Die Fragmente der Vorsokratiker, 2te Aufl., 2 vol., Berlin (Weidmann), 1906-07.

90a J. Dieudonné. Une généralisation des espaces compacts, Journ. de Math., (9), v. XXIII (1944), pp.~65-76.

90b J. Dieudonné. La géométrie des groupes classiques (Erg. der Math. Neue Folge, Heft 5), Berlin-Göttingen-Heidelberg (Springer), 1955.

91a Diophanti Alexandrini Opera Omnia \ldots. 2 vol., ed.~P. Tannery, Lipsiae (Teubner), 1893-95.

91b Diophante d'Alexandrie. Transl. P. Ver Eecke, Bruges (Desclée de Brouwer), 1926.

92a P. G. Lejeune-Dirichlet. Werke, 2 vol., Berlin (Reimer), 1889-1897.

93 J. L. Doob. Probability in function space, Bull. Amer. Math. Soc., v. LIII (1947), pp.~15-30.

94a P. du Bois-Reymond. Sur la grandeur relative des infinis des fonctions, Ann. di Mat. (2), v. IV (1871), pp.~338-353.

94b P. du Bois-Reymond. Üeber asymptotische Werthe, infinitäre Approximationen und infinitäre Auflösung von Gleichungen, Math. Ann., v. VIII (1875), pp.~362-414.

95 N. Dunford. Uniformity in linear spaces, Trans. Amer. Math. Soc., v. XLIV (1938), pp.~305-356.

96 N. Dunford and B. Pettis. Linear Operations on summable functions, Trans. Amer. Math. Soc., v. XLVII (1940), pp.~323-392.

97 W. Dyck. Gruppentheoretische Studien, Math. Ann., v. XX (1882), pp.~1-44.

98a E. Dynkin. Calculation of the coefficients of the Campbell-Hausdorff formula (in Russian), Dokl. Akad. Nauk., v. LVII (1947), pp.~323-326.

98b E. Dynkin. Normed Lie algebras and analytic groups (in Russian), Uspehi Mat. Nauk, v. V (1950), pp.~135-186 (English transl.: Amer. Math. Soc. Transl., (1) vol.~9, pp.~470-534).

99 D. Egoroff. Sur les suites de fonctions mesurables, C. R. Acad. Sci., v. CLII (1911), pp.~244-246.

100 M. Eichler. Quadratische Formen und orthogonale Gruppen, Berlin-Göttingen-Heidelberg (Springer), 1952.

101 A. Einstein. Investigations on the theory of the Brownian movement, New York (Dover), 1956.

102a G. Eisenstein. Beweis der Reciprocitätsgesetze für die cubischen Reste in der Theorie der aus dritten Wurzeln der Einheit zusammengesetzen Zahlen, J. de Crelle, v. XXVII (1844), pp.~289-310.

102b G. Eisenstein. Zur Theorie der quadratischen Zerfällung der Primzahlen 8n + 3, 7n + 2 and 7n + 4, J. de Crelle, v. XXXVII (1848), pp.~97-126.

102c G. Eisenstein. Über einige allgemeine Eigenschaften der Gleichung von welcher die Teilung der ganzen Lemniscate abhängt, nebst Anwendung derselben auf die Zahlentheorie, J. de Crelle, v. XXXIX (1850), pp.~160-179 and 224-287.

103 G. Eneström. Kleine Bemerkung zur letzten Auflage von Cantors Vorlesungen zur Geschichte der Mathematik, Bibl. Math. (3), v. VIII (1907), pp.~412-413.

104a F. Engel. Über die Definitionsgleichung der continuierlichen Transformationsgruppen, Math. Ann. v. XXVII (1886), pp.~1-57.

104b F. Engel. Die Erzeugung der endlichen Transformationen einer projectiven Gruppe durch die infinitesimalen Transformationen der Gruppe, I, Leipziger Ber., v. XLIV (1892), pp.~279-296; II (mit Beiträgen von E. Study), ibid., v. XLV (1893), pp.~659-696.

105a F. Engel und P. Stäckel. Die Theorie der Parallellinien von Euklid bis auf Gauss, Leipzig (Teubner), 1895.

105b F. Engel und P. Stäckel. Urkunden zur Geschichte der nichteuklidischen Geometrien, 2 vol., Leipzig (Teubner), 1898-1913.

106 Enzyklopädie der Mathematischen Wissenschaften. 1te Aufl., 20 vol., Leipzig (Teubner), 1901-1935.

107 Euclidis Elementa. 5 vol., ed.~J. L. Heiberg, Lipsiae (Teubner), 1883-88.

108a L. Euler. Opera Omnia, 46 vol.~appeared (3 series), Leipzig-Berlin-Zürich (Teubner and O. Füssli), 1911-1957.

108b L. Euler. Formulae generales pro translatione quacunque corporum rigidorum, Novi Comm. Acad. Sc. imp. Petrop., v. XX (1776), pp.~189-207.

109 P. Fermat. Oeuvres, 5 vol., Paris (Gauthier-Villars), 1891-1922.

110 X. Fernique. Processus linéaires, processus généralisés, Ann. Inst. Fourier, v. XVII (1967), pp.~1-92.

111 E. Fischer. Sur la convergence en moyenne, C. R. Acad. Sci., v. CXLIV (1907), pp.~1022-1024.

112 J. B. Fourier. Oeuvres, 2 vol., Paris (Gauthier-Villars), 1888-90.

113a A. Fraenkel. Über die Teiler der Null und die Zerlegung von Ringen, J. de Crelle, v. CXLV (1914), pp.~139-176.

113b A. Fraenkel. Zu den Grundlagen der Cantor-Zermeloschen Mengenlehre, Math. Ann., v. LXXXVI (1922), pp.~230-237.

113c A. Fraenkel. Einleitung in die Mengenlehre, 3te Aufl., Berlin (Springer), 1928.

114 D. C. Fraser. Newton's Interpolation Formulas, Journ. Inst. Actuaries, v. LI (1918), pp.~77-106 and pp.~211-232 and v. LVIII (1927), pp.~53-95 (articles reprinted as volume, London, s. d.).

115a M. Fréchet. Sur quelques points du calcul fonctionnel, Rend. Circ. Mat. Palermo, v. XXII (1906), pp.~1-74.

115b M. Fréchet. Les ensembles abstraits et le Calcul fonctionnel, Rend. Circ. Mat. Palermo, v. XXX (1910), pp.~1-26

115c M. Fréchet. Sur l'intégrale d'une fonctionnelle étendue à un ensemble abstrait, Bull. Soc. Math. de France, v. XLIII (1915), pp.~248-265.

115d M. Fréchet. Les familles et fonctions additives d'ensembles abstraits, Fund. Math., v. IV (1923), pp.~229-265 and v. V (1924), pp.~206-251.

116 I. Fredholm. Sur une classe d'équations fonctionnelles, Acta Mathematica, v. XXVII (1903), pp.~365-390.

117a G. Frege. Begriffsschrift eine der arithmetischen nachgebildete Formelsprache des reinen Denkens, Halle, 1879.

117b G. Frege. Die Grundlagen der Arithmetik, 2nd ed.~with an English translation by J. L. Austin, New-York, 1950.

117c G. Frege. Grundgesetze der Arithmetik, begriffschriftlich abgeleitet, 2 vol., Iena, 1893-1903.

118 R. Fricke und F. Klein. Theorie der automorphen Funktionen, Leipzig (Teubner), 1897.

119 G. Frobenius. Gesammelte Abhandlungen (ed.~J. P. Serre), 3 vol., Berlin-Heidelberg-New York (Springer), 1968.

120 G. Frobenius und L. Stickelberger. Ueber Gruppen von vertauschbaren Elementen, J. de Crelle, v. LXXXVI (1879), pp.~217-262.

121 G. Fubini. Sugli integrali multipli, Rendic. Acc. dei Lincei (5), v. XVI (1907), pp.~608-614.

122a Galileo Galilei. Discorsi e Dimostrazioni \ldots, Leyden (Elzevir), 1638.

122b Galileo Galilei. Opere, Ristampa della Ed. Nazionale, 20 vol., Firenze (Barbera), 1929-1939.

123 E. Galois. Écrits et mémoires mathématiques (ed.~R. Bourgne and J. Y. Azra), Paris (Gauthier-Villars), 1962.

124a C. F. Gauss. Werke, 12 vol., Göttingen, 1870-1927.

124b Die vier Gauss'schen Beweise für die Zerlegung ganzer algebraischer Functionen in reelle Factoren ersten oder zweiten Grades (Ostwald's Klassiker, no 14), Leipzig (Teubner), 1904.

125a I. Gelfand. Abstrakte Funktionen und lineare Operatoren, Mat. Sborn. (N.S.), v. IV (1938), pp.~235-284.

125b I. Gelfand. Generalized stochastic processes (in Russian), Dokl. Akad. Nauk, v. C (1955), pp.~853-856.

125c I. Gelfand. Some problems of functional analysis (in Russian), Uspehi Mat. Nauk, v. XI (1956), pp.~3-12 (English transl.: Amer. Math. Soc. Transl., (2), vol.~16 (1960), pp.~315-324).

126 I. Gelfand and N. Ya Vilenkin. Generalized functions, vol.~IV, New York (Academic Press), 1964.

127 G. Gentzen. Die gegenwärtige Lage in der mathematischen Grundlagenforschung. Neue Fassung des Widerspruchsfreiheitsbeweises für die reine Zahlentheorie (Forschungen zur Logik \ldots, Heft 4, Leipzig (Hirzel), 1938).

128 Giorgini. Sopra alcune proprietà de' piani de' momenti \ldots, Mem. Soc. Ital. d.~Sc. res. in Modena, v. XX (1828), pp.~243-254.

129 A. Girard. Invention nouvelle en Algèbre, Amsterdam, 1629 (reprinted ed.~Bierens de Haan, Leyde, 1884).

130a K. Gödel. Über formal unentscheidbare Sätze der Principia Mathematica und verwandter Systeme, Monatsh. für Math. und Phys., v. XXXVIII (1931), pp.~173-198.

130b K. Gödel. The consistency of the axiom of choice and of the generalized continuum hypothesis (Ann. of Math. Studies, no 3), Princeton, 1940.

130c K. Gödel. What is Cantor's continuum hypothesis? Amer. Math. Monthly, v. LIV (1947), pp.~515-525.

131 D. Gorenstein. Finite groups, New York (Harper and Row), 1968.

132 E. Goursat. Sur les substitutions orthogonales et les divisions régulières de l'espace, Ann. Ec. Norm. Sup., (3), v. VI (1889), pp.~9-102.

133 J. P. Gram. Ueber die Entwicklung reeller Functionen in Reihen mittelst der Methode der kleinsten Quadrate, J. de Crelle, v. XCIV (1883), pp.~41-73.

134 H. Grassmann. Gesammelte Werke, 3 vol., Leipzig (Teubner), 1894-1911.

135 P. Gregorii a Sancto Vicentio. Opus Geometricum Quadraturae Circuli et Sectionum Coni \ldots, 2 vol., Antverpiae, 1647.

136a J. Gregory. Vera Circuli et Hyperbolae Quadraturae \ldots, Pataviae, 1667.

136b J. Gregory. Geometriae Pars Universalis, Pataviae, 1668.

136c J. Gregory. Exercitationes Geometricae, London, 1668.

136d James Gregory Tercentenary Memorial Volume, containing his correspondence with John Collins and his hitherto unpublished mathematical manuscripts \ldots. Ed. H. W. Turnbull, London (Bell and sons), 1939.

137 H. Grell. Beziehungen zwischen den Idealen verschiedener Ringe, Math. Ann., v. XCVII (1927), pp.~490-523.

138a A. Grothendieck. Éléments de géométrie algébrique, I, Publ. math. I. H. E. S., no 4 (1960).

138b A. Grothendieck. Produits tensoriels topologiques et espaces nucléaires, Mem. Amer. Math. Soc., no 16 (1955).

138c A. Grothendieck. Espaces vectoriels topologiques, 3rd ed., São Paulo (Publ. Soc. Mat. São Paulo), 1964.

139 A. Haar. Der Maassbegriff in der Theorie der kontinuierlichen Gruppen, Ann. of Math., v. XXXIV (1933), pp.~147-169.

140 J. Hachette and S. Poisson. Addition au mémoire précédent, Journ. de l'Ec. Polytechn., cahier 11 (an X), pp.~170-172.

141 J. Hadamard. Sur certaines applications possibles de la théorie des ensembles, Verhandl. Intern. Math. Kongress, Zürich, 1898, pp.~201-202.

142 H. Hahn. Ueber lineare Gleichungssysteme in linearen Räumen, J. de Crelle, v. CLVII (1927), pp.~214-229.

143 M. Hall. A basis for free Lie rings and higher commutators in free groups, Proc. Amer. Math. Soc., v. I (1950), pp.~575-581.

144 P. Hall. A contribution to the theory of groups of prime power order, Proc. Lond. Math. Soc., (3), v. IV (1932), pp.~29-95.

145a W. R. Hamilton. Lectures on quaternions, Dublin, 1853.

145b W. R. Hamilton. Memorandum respecting a new system of roots of unity, Phil. Mag., (4), v. XII (1856), p.~446.

146 H. Hankel. Theorie der complexen Zahlensysteme, Leipzig (Voss), 1867.

147 G. H. Hardy. Orders of infinity (Cambridge tracts, no 12), 2nd ed., Cambridge University Press, 1924.

148 G. H. Hardy and E. M. Wright. An introduction to the Theory of Numbers, Oxford, 1938.

149a Harish-Chandra. On some applications of the universal enveloping algebra of a semi-simple Lie algebra, Trans. Amer. Math. Soc., v. LXX (1951), pp.~28-96.

149b Harish-Chandra. On a lemma of Bruhat, Journ. de Math., (9), v. XXXV (1956), pp.~203-210.

150a H. Hasse. Ueber die Darstellbarkeit von Zahlen durch quadratischen Formen im Körper der rationalen Zahlen, J. de Crelle, v. CLII (1923), pp.~129-148.

150b H. Hasse. Ueber die Äquivalenz quadratischer Formen in Körper der rationalen Zahlen, J. de Crelle, v. CLII (1923), pp.~205-224.

150c H. Hasse. Kurt Hensels entscheidender Anstoss zur Entdeckung des Lokal-Global-Prinzips, J. de Crelle, v. CCIX (1960), pp.~3-4.

150d H. Hasse. Zahlentheorie, Berlin (Akad. Verlag), 1949.

151 H. Hasse und H. Scholz. Die Grundlagenkrise der grieschischen Mathematik, Charlottenburg (Pan-Verlag), 1928 (= Kant-Studien, v. XXXIII (1928), pp.~4-72).

152a F. Hausdorff. Die symbolische Exponentialformal in der Gruppentheorie, Leipziger Ber., v. LVIII (1906), pp.~19-48.

152b F. Hausdorff. Grundzüge der Mengenlehre, Leipzig (Veit), 1914.

152c F. Hausdorff. Mengenlehre, Berlin (de Gruyter), 1927.

153a T. Heath. A History of Greek Mathematics, 2 vol., Oxford 1921.

153b T. Heath. Apollonius of Perga, Treatise on conic sections, Cambridge University Press, 1896.

153c T. Heath. The method of Archimedes, Cambridge, 1912.

153d T. Heath. Mathematics in Aristotle, Oxford (Clarendon Press), 1949.

153e T. Heath. The thirteen books of Euclid's Elements \ldots, 3 vol., Cambridge, 1908.

153f T. Heath. Diophantus of Alexandria, 2nd ed., Cambridge, 1910.

154a E. Heine. Ueber trigonometrische Reihen, J. de Crelle, v. LXXI (1870), pp.~353-365.

154b E. Heine. Die Elemente der Functionenlehre, J. de Crelle, v. LXXIV (1872), pp.~172-188.

155 G. Heinrich. James Gregorys ``Vera circuli et hyperbolae quadratura'', Bibl. Math., (3), v. II (1901), pp.~77-85.

156 E. Helly. Ueber Systeme linearer Gleichungen mit unendlich vielen Unbekannten, Monatsh. für Math. und Phys., v. XXXI (1921), pp.~60-91.

157a K. Hensel. Über eine neue Begründung der Theorie der algebraischen Zahlen, Jahresber. der D. M. V., v. VI (1899), pp.~83-88.

157b K. Hensel. Ueber die Fundamentalgleichung und die ausserwesentlichen Diskriminantentheiler eines algebraischen Körpers, Gött. Nachr. (1897), pp.~254-260.

157c K. Hensel. Neue Grundlagen der Arithmetik, J. de Crelle, v. CXXVII (1902), pp.~51-84.

157d K. Hensel. Über die arithmetische Eigenschaften der algebraischen und transzendenten Zahlen, Jahresber. der D. M. V., v. XIV (1905), pp.~545-558.

157e K. Hensel. Ueber die arithmetischen Eigenschaften Zahlen, Jahresber. der D. M. V., v. XVI (1907), pp.~299-319, 388-393, 474-496.

157f K. Hensel. Theorie der algebraischen Zahlen, Leipzig (Teubner), 1908.

158 J. Herbrand. Recherches sur la théorie de la démonstration, Trav. Soc. Sci. Lett. Varsovie, cl. II (1930), pp.~33-160.

159 C. Hermite. Oeuvres, 4 vol., Paris (Gauthier-Villars), 1905-1917.

160 C. Hermite, T. Stieltjes. Correspondance, 2 vol., Paris (Gauthier-Villars), 1905.

161 J. Hessel. Krystallometrie oder Krystallonomie und Krystallographie (1830, repr. in Ostwald's Klassiker, nos 88 and 89, Leipzig (Teubner), 1897).

162 A. Heyting. Mathematische Grundlagenforschung. Intuitionismus. Beweistheorie (Erg. der Math., Bd. 3), Berlin (Springer), 1934.

163a D. Hilbert. Gesammelte Abhandlungen, 3 vol., Berlin (Springer), 1932-35.

163b D. Hilbert. Grundzüge einer allgemeinen Theorie der Integralgleichungen, 2nd ed., Leipzig-Berlin (Teubner), 1924.

163c D. Hilbert. Grundlagen der Geometrie, 7th ed., Leipzig-Berlin (Teubner), 1930.

164 D. Hilbert und W. Ackermann. Grundzüge der theoretischen Logik, 3te Aufl., Berlin (Springer), 1949.

165 O. Hölder. Zurückführung einer beliebigen algebraischen Gleichung auf eine Kette von Gleichungen, Math. Ann., v. XXXIV (1889), pp.~26-56.

166 R. Hooke. Linear \(p\) -adic groups and their Lie algebras, Ann. of Math., v. XLIII (1942), pp.~641-655.

167 C. Hopkins. Rings with minimal conditions for left ideals, Ann. of Math., (2), v. XL (1939), pp.~712-730.

168 A. Hurwitz. Über die Erzeugung der Invarianten durch Integration, Gött. Nachr., 1897, pp.~71-90 (= Math. Werke, v. II, pp.~546-564).

169a Christiani Hugenii, Zulichemii Philosophi vere magni, Dum viveret Zelemii Toparchae, Opera \ldots. 4 tomes in 1 vol., Lugd. Batav., 1751.

169b C. Huygens. Oeuvres complètes, 22 vol., The Hague (M. Nijhoff), 1888-1950.

170 N. Iwahori and H. Matsumoto. On some Bruhat decompositions and the structure of the Hecke ring of \(p\) -adic Chevalley groups, Publ. math. I. H. E. S., no 25 (1965), pp.~5-48.

171 C. G. J. Jacobi. Gesammelte Werke, 7 vol., Berlin (G. Reimer), 1881-91.

172a N. Jacobson. Rational methods in the theory of Lie algebras, Ann. of Math., v. XXXVI (1935), pp.~875-881.

172b N. Jacobson. Classes of restricted Lie algebras of characteristic p, II, Duke Math. Journ., v. X (1943), pp.~107-121.

172c N. Jacobson. A topology for the set of primitive ideals in an arbitrary ring, Proc. Nat. Acad. Sci. U. S. A., v. XXXI (1945), pp.~333-338.

172d N. Jacobson. Structure of rings (Amer. Math. Soc. Coll. Public., v. 37), Providence, 1956.

173 B. Jessen. The theory of integration in a space of an infinite number of dimensions, Acta Math., v. LXIII (1934), pp.~249-323.

174a C. Jordan. Traité des substitutions et des equations algébriques, 2nd ed., Paris (Gauthier-Villars et A. Blanchard), 1957.

174b C. Jordan. Mémoire sur les groupes de mouvements, Ann. di Mat., (2), v. II (1868-69), pp.~167-215 and 322-345 (= Oeuvres, v. IV, pp.~231-302, Paris (Gauthier-Villars 4), 1964.)

174c C. Jordan. Cours d'Analyse de l'École Polytechnique, 3rd ed., 3 vol., Paris (Gauthier-Villars), 1909-15.

175 H. Jung. Algebraischen Flächen, Hannover (Helwing), 1925.

176 J. P. Kahane. Séries de Fourier aléatoires, Sém. Bourbaki, no. 200, 12th year, 1959-60, New York (Benjamin).

177 S. Kakutani. Notes on infinite product measure spaces, II, Proc. Imp. Acad. Japan, v. XIX (1943), pp.~184-188.

178 I. Kant. Werke, ed.~E. Cassirer, 11 vol., Berlin (B. Cassirer), 1912-1921.

179a J. Kepler. Stereometria Doliorum, 1615.

179b J. Kepler. Neue Stereometrie der Fäser (Ostwald's Klassiker, no. 165), Leipzig (Engelmann), 1908.

180 W. Killing. Die Zusammensetzung der stetigen endlichen Transformationsgruppen : I) Math. Ann., v. XXXI (1888), pp.~252-290; II) ibid., v. XXXIII (1889), pp.~1-48; III) ibid., v. XXXIV (1889), pp.~57-122; IV) ibid., v. XXXVI (1890), pp.~161-189.

181 S. Kleene. Introduction to metamathematics, New York, 1952.

182 F. Klein. Gesammelte mathematische Abhandlungen, 3 vol., Berlin (Springer), 1921-23.

183 A. Kolmogoroff. Grundgebriffe der Wahrscheinlichkeitsrechnung (Erg. der Math., Bd. 2), Berlin (Springer), 1933.

184 G. Köthe. Neubegründung der Theorie der vollkommenen Räume, Math. Nachr., v. IV (1951), pp.~70-80.

185 E. Kötter. Die Entwicklung der synthetischen Geometrie, Leipzig (Teubner), 1901 (= Jahresbericht der D. M. V., v. V, 2tes Heft).

186a L. Kronecker. Werke, 5 vol., Leipzig (Teubner), 1895-1930.

186b L. Kronecker. Vorlesungen über die Theorie der Determinanten \ldots, Leipzig (Teubner), 1903.

187a W. Krull. Über verallgemeinerte endliche Abelsche Gruppen, Math. Zeitschr., v. XXIII (1925), pp.~161-196.

187b W. Krull. Theorie und Anwendung der verallgemeinerten Abelschen Gruppen, Sitzungsber. Heidelberger Akad. Wiss., 1926, no. 1, 32 pp.

187c W. Krull. Zur Theorie der allgemeinen Zahlringe, Math. Ann., v. XCIX (1928), pp.~51-70.

187d W. Krull. Galoische Theorie der unendlichen algebraischen Erweiterungen, Math. Ann., v. C (1928), pp.~687-698.

187e W. Krull. Primidealketten in allgemeine Ringbereichen, Sitzungsber. Heidelberg Akad. Wiss., 1928.

187f W. Krull. Allgemeine Bewertungstheorie, J. de Crelle, v. CLXVII (1931), pp.~160-196.

187g W. Krull. Beiträge zur Arithmetik kommutativer Integritätsbereiche III, Math. Zeitschr., v. XLII (1937), pp.~745-766.

187h W. Krull. Dimensionstheorie in Stellenringe, J. de Crelle, v. CLXXIX (1938), pp.~204-226.

187i W. Krull. Idealtheorie (Erg. der Math., Bd 4), Berlin (Springer), 1935.

188a E. Kummer. Sur les nombres complexes qui sont formés avec les nombres entiers réels et les racines de l'unité, Journ. de Math., (1), v. XII (1847), pp.~185-212.

188b E. Kummer. Zur Theorie der complexen Zahlen, J. de Crelle, v. XXXV (1847), pp.~319-326.

188c E. Kummer. Ueber die Zerlegung der aus Wurzeln der Einheit gebildeten complexen Zahlen in Primfactoren, J. de Crelle, v. XXXV (1847), pp.~327-367.

188d E. Kummer. Mémoire sur les nombres complexes composés de racines de l'unité et des nombres entiers, Journ. de Math., (1), v. XVI (1851), pp.~377-498.

188e E. Kummer. Über die allgemeinen Reciprocitätsgesetze unter den Resten und Nichtresten der Potenzen deren Grad eine Primzahl ist, Abhandl. der Kön. Akad. der Wiss. zu Berlin (1859), Math. Abhandl., pp.~19-159.

189a K. Kuratowski. Une méthode d'élimination des nombres transfinis des raisonnements mathématiques, Fund. Math., v. V (1922), pp.~76-108.

189b K. Kuratowski. Topologie, I, 2nd ed., Warszawa-Vroclaw, 1948.

190 J. Kürschak. Über Limesbildung und allgemeine Körpertheorie, J. de Crelle, v. CXLII (1913), pp.~211-253.

191 J. L. Lagrange. Oeuvres, 14 vol., Paris (Gauthier-Villars), 1867-1892.

192 E. Laguerre. Oeuvres, 2 vol., Paris (Gauthier-Villars), 1898-1905.

193 P. S. Laplace. Oeuvres, 14 vol., Paris (Gauthier-Villars), 1878-1912.

194 E. Lasker. Zur Theorie der Moduln und Ideale, Math. Ann., v. LX (1905), pp.~20-116.

195a M. Lazard. Quelques calculs concernant la formule de Hausdorff, Bull. Soc. Math. de France, v. XCI (1963), pp.~435-451.

195b M. Lazard. Groupes analytiques p-adiques, Publ. Math. I. H. E. S., no. 26 (1965), pp.~389-603.

196a H. Lebesgue. Intégrale, longueur, aire, Ann. di Mat. (3), v. VII (1902), pp.~231-259.

196b H. Lebesgue. Sur les séries trigonométriques, Ann. Ec. Norm. Sup., (3), v. XX (1903), pp.~453-485.

196c H. Lebesgue. Leçons sur l'Intégration et la recherche des fonctions primitives, Paris (Gauthier-Villars), 1904.

196d H. Lebesgue. Sur le problème de Dirichlet, Rend. Circ. Mat. Palermo, v. XXIV (1907), pp.~371-402.

196e H. Lebesgue. Sur l'intégration des fonctions discontinues, Ann. Ec. Norm. Sup. (3), v. XXVII (1910), pp.~361-450.

197 L. Le Cam. Convergence in distribution of stochastic processes, Univ. Calif. Publ. Statistics, no. 11 (1957), pp.~207-236.

198a G. W. Leibniz. Mathematische Schriften, 7 vol., ed.~C. I. Gerhardt, Berlin-Halle (Ascher-Schmidt), 1849-63.

198b G. W. Leibniz. Philosophische Schriften, 7 vol., ed.~C. I. Gerhardt, Berlin, 1840-90.

198c G. W. Leibniz. Opuscules et fragments inédits, ed.~L. Couturat, Paris (Alcan), 1903.

198d Der Briefwechsel von Gottfried Wilhelm Leibniz mit Mathematikern. V. I, herausgegeben von C. I. Gerhardt, Berlin (Mayer und Müller), 1899.

199 B. Levi. Intorno alla teoria degli aggregati, R. Ist. Lombardo Sci. Lett. Rendic. (2), v. XXXV (1902), pp.~863-868.

200 E. E. Levi. Sulla struttura dei Gruppi finiti e continui, Atti Accad. Sci. Torino, v. XL (1905), pp.~551-565 (= Opere, v. I, p.~101-115).

201a P. Lévy. Leçons d'Analyse fonctionnelle, Paris (Gauthier-Villars), 1922) (2dn ed., under the title: Problèmes concrets d'Analyse fonctionnelle, 1951).

201b P. Lévy. Processus stochastiques et mouvement brownien, Paris (Gauthier-Villars), 1948.

201c P. Lévy. Le mouvement brownien, Mémor. des Sci. Math., v. CXXVI (1954), Paris (Gauthier-Villars).

202 S. Lie. Gesammelte Abhandlungen, 7 vol., Leipzig (Teubner).

203 S. Lie und F. Engel. Theorie der Transformationsgruppen, 3 vol., Leipzig (Teubner), 1888-1893.

203bis J. Lindenstrauss und L. Tzafriri. Classical Banach spaces, v. I, Berlin-Heidelberg-New York (Springer), 1977.

204a J. Liouville. Sur le développement des fonctions ou parties de fonctions en séries dont les divers termes sont assujettis à satisfaire à une même équation différentielle du second ordre contenant un paramètre variable, Journ. de Math. (1), v. I (1836), pp.~253-265 and v. II (1837), pp.~16-35 and 418-436.

204b J. Liouville. D'un théorème dû à M. Sturm et relatif à une classe de fonctions transcendants, Journ. de Math. (1), v. I (1836), pp.~269-277.

204c J. Liouville. Sur des classes très étendues de quantités dont la valeur n'est ni algébrique ni même réductible à des irrationnelles algébriques, Journ. de Math. (1), v. XVI (1851), pp.~133-142.

205a R. Lipschitz. De explicatione per serias trigonometricas instituenda functionum unius variabilis arbitrariarum, J. de Crelle, v. LXIII (1864), pp.~296-308 (French transl. by P. Montel Acta Math., v. XXXVI (1912), pp.~261-295).

205b R. Lipschitz. Untersuchungen ueber die Summen von Quadraten, Bonn, 1886 (summarized in French in Bull. Sci. Math., v. XXI (1886), pp.~163-183).

206 N. Lobatschevsky. Pangeometrie (Ostwald's Klassiker, no. 130), Leipzig (Engelmann), 1902.

207 L. H. Loomis. An introduction to abstract harmonic analysis, London-New York-Toronto (van Nostrand), 1953.

208 G. Loria. Le ricerche inedite de Evangelista Torricelli sopra la curve logaritmica, Bibl. Math. (3), v. I (1900), pp.~75-89.

209 N. Lusin. Leçons sur les ensembles analytiques et leurs applications, Paris (Gauthier-Villars), 1930.

210 E. Lutz. Sur l'équation \(y^{2} = x^{3} - Ax - B\) dans les corps \(p\) -adiques, J. de Crelle, v. CLXXVII (1937), pp.~237-247.

211 F. S. Macaulay. On the resolution of a given modular system into primary systems including some properties of Hilbert numbers, Math. Ann., v. LXXIV (1913), pp.~66-121.

212 A. Macbeath and S. Swierczkowski. Limits of lattices in a compactly generated group, Can. Journ. of Math., v. XII (1960), pp.~427-437.

213a G. W. Mackey. On infinite-dimensional spaces, Trans. Amer. Math. Soc., v. LVII (1945), pp.~155-207.

213b G. W. Mackey. On convex topological spaces, Trans. Amer. Math. Soc., v. LX (1946), pp.~519-537.

214 C. Maclaurin. A treatise of fluxions, Edinburgh, 2 vol., 1742.

215a W. Magnus. Beziehungen zwischen Gruppen und Idealen in einen speziellen Ring, Math. Ann., v. CXI (1935), pp.~259-280.

215b W. Magnus. Über Beziehungen zwischen höheren Kommutatoren, J. de Crelle, v. CLXXVII (1937), pp.~105-115.

216 W. Magnus, A. Karass and D. Solitar. Combinatorial group theory, New York (Interscience), 1966.

217 D. Mahnke. Neue Einblicke in die Entdeckungsgeschichte der höheren Analysis, Abh. Preuss. Akad. der Wiss. Phys.-Math. Klasse, 1925, no. 1, 64 pp., Berlin, 1926.

218 F. Masères. Scriptores Logaritmici, 6 vol., London, 1791-1807.

219 L. Maurer. Über allgemeinere Invarianten-Systeme, Sitzungsber. München, v. XVIII (1888), pp.~103-150.

220 N. Mercator. Logarithmotechnia\ldots{} cui nunc accedit vera quadratura hyperbolae \ldots, Londini, 1668.

221a H. Minkowski. Gesammelte Abhandlungen, 2 vol., Leipzig-Berlin (Teubner), 1911.

221b H. Minkowski. Geometrie der Zahlen, Leipzig (Teubner), 1896.

222 R. A. Minlos. Generalized random processes and their extension to a measure (in Russian), Trudy Mosk. Mat. Obschtsch., v. VIII (1959), pp.~497-518 (= Selected transl. in math. statistics and probability, v. III (19), pp.~291-313, Amer. Math. Soc.).

223 A. F. Möbius. Gesammelte Werke, 4 vol., Leipzig (Hirzel), 1885-87.

224a T. Molien. Ueber Systeme höherer complexer Zahlen, Math. Ann., v. XLI (1893), pp.~83-156.

224b T. Molien. Über die Invarianten der linearen Substitutionsgruppen, Berliner Sitzungsber., 1897, pp.~1152-1156.

225 G. Monge. Géométrie descriptive, Paris, 1798 (new ed.~Paris (Gauthier-Villars), 1922).

226 D. Montgomery and L. Zippin. Topological transformation groups, New York (Interscience), 1955.

227 E. H. Moore and H. L. Smith. A general theory of limits, Amer. Journ. of Math., v. XLIV (1922), pp.~102-121.

228 S. G. Morley. The ancient Maya, Stanford University Press, 1946.

229 E. Nelson. Dynamical theories of brownian motion, Mathematical Notes, Princeton (University Press), 1967.

230a J. Napier. Mirifici logarithmorum canonis descriptio, Lyon, 1619 (reproduced in {[}218{]}, v. VI, pp.~475-623).

230b Napier Tercentenary Memorial Volume. London, 1915.

231 E. Netto. Zur Theorie der Elimination, Acta Math., v. VII (1885), pp.~101-104.

232 O. Neugebauer. Vorlesungen über die Geschichte der antiken Mathematik, Bd. I: Vorgrieschische Mathematik, Berlin (Springer), 1934.

233a I. Newton. Opuscula, 3 vol., Lausanne-Genève (M. Bousquet), 1744.

233b I. Newton. Philosophiae Naturalis Principia Mathematica, London, 1687 (new ed., Glasgow, 1871).

233c I. Newton. Mathematical principles of natural philosophy, transl. into English by A. Motte in 1729, Univ. of California, 1946.

233d The mathematical papers of Isaac Newton. Vol. I: 1664-1666 (ed.~D. Whiteside), Cambridge (Univ. Press), 1967.

234 J. Nielsen. Die Isomorphismengruppen der freien Gruppen, Math. Ann., v. XCI (1924), pp.~169-209.

235 O. Nikodym. Sur une généralisation des intégrales de M. Radon, Fund. Math., v. XV (1930), pp.~131-179.

236a E. Noether. Idealtheorie in Ringbereichen, Math. Ann., v. LXXXIII (1921), pp.~24-66.

236b E. Noether. Abstrakter Aufbau der Idealtheorie in algebraischen Zahl- und Funktionenkörper, Math. Ann., v. XCVI (1926), pp.~26-61.

236c E. Noether. Hyperkomplexe Grössen und Darstellungstheorie, Math. Zeitschr., v. XXX (1929), pp.~641-692.

236d E. Noether. Nichtkommutative Algebra, Math. Zeitschr., v. XXXVII (1933), pp.~514-541.

237 E. Noether und R. Brauer. Über minimale Zerfallungskörper irreduzibler Darstellungen, Berliner Sitzungsber., 1927, pp.~221-228.

238 E. Noether und W. Schmeidler. Moduln in nichtkommutativen Bereichen, Math. Zeitschr., v. VIII (1920), pp.~1-35.

239 M. Noether. Über einen Satz aus der Theorie der algebraischen Funktionen, Math. Ann., v. VI (1873), pp.~351-359.

240 W. Osgood. Non uniform convergence and the integration of series term by term, Amer. Journ. of Math., v. XIX (1897), pp.~155-190.

241 P. Osmond. Isaac Barrow. His life and time, London, 1944.

242 A. Ostrowski. Über einige Lösungen der Funktionalgleichung \(\phi(x)\phi(y) = \phi(x.y)\) , Acta Math., v. XLI (1917), pp.~271-284.

243 R. E. A. C. Paley and N. Wiener. Fourier transforms in the complex domain, Amer. Math. Soc. Coll. Publ., no. 19, New York, 1934.

244 B. Pascal. Oeuvres, 14 vol., ed.~Brunschvicg, Paris (Hachette), 1904-14.

245 M. Pasch und M. Dehn. Vorlesungen über neuere Geometrie, 2te Aufl., Berlin (Springer), 1926.

246a G. Peano. Applicazioni geometriche del calcolo infinitesimale, Torino, 1887.

246b G. Peano. Calcolo geometrico secondo l'Ausdehnungslehre di Grassmann, preceduto dalle operazioni della logica deduttiva, Torino, 1888.

246c G. Peano. Arithmeticas principia, novo methodo exposita, Torino, 1889.

246d G. Peano. I principii di Geometria, logicamente expositi, Torino, 1889.

246e G. Peano. Démonstration de l'intégrabilité des équations différentielles ordinaires, Math. Ann., v. XXXVII (1890), pp.~182-228.

246f G. Peano. Formulaire de Mathématiques, 5 vol., Torino, 1895-1905.

247 B. Peirce. Linear associative algebra, Amer. Journ. of Math., v. IV (1881), pp.~97-221.

248a C. S. Peirce. Upon the logic of mathematics, Proc. Amer. Acad. of Arts and Sci., v. VII (1865-68), pp.~402-412.

248b C. S. Peirce. On the algebra of logic, Amer. Journ. of Math., v. III (1880), pp.~49-57.

248c C. S. Peirce. On the relative forms of the algebras, Amer. Journ. of Math., v. IV (1881), pp.~221-225.

248d C. S. Peirce. On the algebras in which division is unambiguous, Amer. Journ. of Math., v. IV (1881), pp.~225-229.

248e C. S. Peirce. On the algebra of logic, Amer. Journ. of Math., v. VII (1884), pp.~190-202.

249 J. Pfanzagl and W. Pierlo. Compact systems of sets, Lecture Notes in Math., no. 16 (1966), Berlin (Springer).

250 Plato. La République, transl. E. Chambry, 2 vol., Paris (Les Belles Lettres), 1932-49.

251a H. Poincaré. Oeuvres, 11 vol., Paris (Gauthier-Villars), 1916-1956.

251b H. Poincaré. Les méthodes nouvelles de la mécanique céleste, 3 vol., Paris (Gauthier-Villars), 1893-1899.

251c H. Poincaré. Science et hypothèse, Paris (Flammarion), 1902.

251d H. Poincaré. La valeur de la Science, Paris (Flammarion), 1905.

251e H. Poincaré. Science et méthode, Paris (Flammarion), 1908.

252 J. V. Poncelet. Traité des propriétés projectives des figures, 2 vol., 2nd ed., Paris (Gauthier-Villars), 1865.

253 L. S. Pontrjagin. Topological groups, Princeton (Univ. Press), 1939.

254 Ju. V. Prokhorov. Convergence of random processes and limit theorems in probability theory, Theor. Probab. Appl., v. I (1956), pp.~156-214.

255 Ptolemaei Cl. Opera. Ed. J. L. Heiberg, 2 vol., Lipsiae (Teubner), 1898-1903 (transl. Halma, reprinted., 2 vol., Paris (Hermann), 1927).

256 J. Radon. Theorie und Anwendungen der absolut additiven Mengenfunctionen, Sitzungsber. der math. naturwiss. Klasse der Akad. der Wiss. (Wien), v. CXXII, Abt. II a (1913), pp.~1295-1438.

257 G. Ricci et T. Levi-Civita. Méthodes de calcul différentiel absolu et leurs applications. Math. Ann., v. LIV (1901), pp.~125-201.

258 J. Richard. Les principes des Mathématiques et le problème des ensembles, Riv. Gén. des Sci. pures et appl., v. XVI (1905), pp.~541-543.

259a B. Riemann. Gesammelte mathematische Werke, 2nd ed., Leipzig (Teubner), 1892.

259b B. Riemann. Gesammelte Werke, Nachträge, Leipzig (Teubner), 1902.

259c B. Riemann. in Lettere di E. Betti a P. Tardy, Rend. Accad. dei Lincei (5), v. XXIV \(^{1}\) (1915), pp.~517-519.

260a F. Riesz. Sur les systèmes orthogonaux de fonctions, C. R. Acad. Sci., v. CXLIV (1907), pp.~615-619.

260b F. Riesz. Stetigkeitsbegriff und abstrakte Mengenlehre, Atti del IV Congresso Intern. dei Matem., Roma, 1908, v. II, pp.~18-24.

260c F. Riesz. Untersuchungen über Systeme integrierbarer Funktionen, Math. Ann., v. LXIX (1910), pp.~449-497.

260d F. Riesz. Sur certains systèmes singuliers d'équations intégrales, Ann. Ec. Norm. Sup. (3), v. XXVIII (1911), pp.~33-62.

260e F. Riesz. Les systèmes d'équations linéaires à une infinité d'inconnues, Paris (Gauthier-Villars), 1913.

260f F. Riesz. Ueber lineare Funktionalgleichungen, Acta Math., v. XLI (1918), pp.~71-98.

260g F. Riesz. Zur Theorie des Hilbertschen Raumes, Acta litt. ac. scient. (Szeged), v. VII (1934-35), pp.~34-38.

260h F. Riesz. Sur quelques notions fondamentales dans la théorie générale des opérations linéaires, Ann. of Math. (2), v. XLI (1940), pp.~174-206.

261 M. Riesz. Sur le problème des moments, 3, Ark. für Math., v. XVII (1922-23), no. 16, 52 pp.

262 S. P. Rigaud. Correspondence of scientific men \ldots, 2 vol., Oxford, 1841-42.

263 G. de Roberval. Ouvrages de Mathématique (Mémoires de l'Académie Royale des Sciences, v. VI, Paris (1730), pp.~1-478).

264 R. Robinson. Plato's consciousness of fallacy, Mind, v. LI (1942), pp.~97-114.

265 P. Ruffini. Opere Matematiche, 3 vol., Ed. Cremonese (Roma), 1953-54.

266 B. Russell and A. N. Whitehead. Principia Mathematica, 3 vol., Cambridge, 1910-13.

267 A. Rustow. Der Lügner, Diss. Erlangen, 1910.

268 S. Saks. Theory of the integral, 2nd ed., New York (Stechert), 1937.

269 P. Alfonso Antonio de Sarasa. Solutio problematis \ldots, Antverpiae, 1649.

270 P. Samuel. La notion de multiplicité en Algèbre et en Géométrie algébrique, Journ. de Math. (9), v. XXX (1951), pp.~159-274.

271 G. Scheffers. Zurückführung complexer Zahlensysteme auf typische Formen, Math. Ann., v. XXXIX (1891), pp.~293-390.

272 E. Schering. Gesammelte mathematische Werke, 2 vol., Berlin (Mayer und Müller), 1902-1909.

273 L. Schäfli. Gesammelte mathematische Abhandlungen, 3 vol., Basel (Birkhäuser), 1950-56.

274a E. Schmidt. Zur Theorie der linearen und nichtlinearen Integralgleichungen. I. Teil: Entwicklung willkürlicher Funktionen nach Systeme vorgeschriebener, Math. Ann., v. LXIII (1907), pp.~433-476.

274b E. Schmidt. Ueber die Auflösung linearer Gleichungen mit unendlich vielen Unbekannten, Rend. Palermo, v. XXV (1908), pp.~53-77.

275a A. Schoenfleis. Entwicklung der Mengenlehre und ihrer Anwendungen, 2nd ed., Leipzig-Berlin (Teubner), 1913.

275b A. Schoenfleis. Die Krisis in Cantor's mathematischem Schaffen, Acta Math., v. L (1927), pp.~1-23.

276a O. Schreier. Abstrakte kontinuierliche Gruppen, Abh. math. Sem. Univ. Hambrug, v. IV (1926), pp.~15-32.

276b O. Schreier. Die Verwandschaft stetiger Gruppen in grossen, Abh. math. Sem. Univ. Hamburg, v. V (1927), pp.~233-244.

277 E. Schröder. Vorlesungen über die Algebra der Logik, 3 vol., Leipzig (Teubner), 1890.

278a F. Schur. Zur Theorie der aus Haupteinheiten gebildeten Komplexen, Math. Ann., v. XXXIII (1889), pp.~49-60.

278b F. Schur. Neue Begründung der Theorie der endlichen Transformationsgruppen, Math. Ann., v. XXXV (1890), pp.~161-197.

278c F. Schur. Zur Theorie der endlichen Transformationsgruppen, Math. Ann., v. XXXVIII (1891), pp.~273-286.

278d F. Schur. Über den analytischen Character der eine endliche continuierliche Transformationsgruppen darstellende Funktionen, Math. Ann., v. XLI (1893), pp.~509-538.

279a I. Schur. Über eine Klasse von Matrices, die sich einer gegebenen Matrix zuordnen lassen, Diss. Berlin, 1901.

279b I. Schur. Über die Darstellung der endlichen Gruppen durch gebrochene lineare Substitutionen, J. de Crelle, v. CXXVII (1904), pp.~20-50.

279c I. Schur. Neue Begründung der Theorie der Gruppencharaktere, Berliner Sitzungsber., 1905, pp.~406-432.

279d I. Schur. Arithmetische Untersuchungen über endliche Gruppen linearer Substitution, Berliner Sitzungsber., 1906, pp.~164-184.

279e I. Schur. Neue Anwendungen der Integralrechnung auf die Probleme der Invariantentheorie, Berliner Sitzungsber., 1924, pp.~189-208, 297-321, 346-355.

280 L. Schwartz. Théorie des distributions (Actual. Scient. et Ind., nos 1091 and 1122), Paris (Hermann), 1950-51.

281 I. Segal. Distributions in Hilbert space and canonical systems of operators, Trans. Amer. Math. Soc., v. LXXXVIII (1958), pp.~12-41.

282 E. Selling. Ueber die idealen Primfactoren der complexen Zahlen, welche aus der Wurzeln einer beliebigen irreductiblen Gleichung rational gebildet sind, Zeitschr. für Math. und Phys., v. X (1865), pp.~17-47.

283a J.-P. Serre. Faisceaux algébriques cohérents, Ann. of Math., v. LXI (1955), pp.~197-278.

283b J.-P. Serre. Géométrie algébrique et géométrie analytique, Ann. Inst. Fourier, v. VI (1956), pp.~1-42.

284 J. A. Serret. Cours d'Algèbre Supérieure, 3rd ed., Paris (Gauthier-Villars), 1866.

285 C. L. Siegel. Symplectic Geometry, Amer. Journ. of Math., v. LXV (1943), pp.~1-86.

286a T. Skolem. Einige Bemerkungen zur axiomatischen Begründung der Mengenlehre, Wiss. Vorträge, 5. Kongress Skand. Math., Helsingfors, 1922, pp.~217-232.

286b T. Skolem. Zur Theorie der associativen Zahlensysteme, Skr. norske Vid. Akad., Oslo, 1927, no. 12, 50 pp.

287 H. J. Smith. Collected Mathematical Papers, 2 vol., Oxford, 1894.

288 Sommer. Introduction à la théorie des nombres algébriques (transl. A. Lévy), Paris (Hermann), 1911.

289 M. Souslin. Sur une définition des ensembles mesurables B sans nombres transfinis, C. R. Acad. Sci., v. CLXIV (1917), pp.~88-91.

290 E. Sparre-Andersen and B. Jessen. On the introduction of measures in infinite product sets, Dansk. Vid. Selbskab. Mat. Fys. Medd., v. XXV (1948), no. 4, pp.~1-7.

291 A. Speiser. Theorie der Gruppen von endlicher Ordnung, 4th ed., Basel (Birkhäuser), 1956.

292 R. Steinberg. Finite reflection groups, Trans. Amer. Math. Soc., v. XCI (1959), pp.~493-504.

293 H. Steinhaus. Les probabilités dénombrables et leur rapport à la théorie de la mesure, Fund. Math., v. IV (1923), pp.~286-310.

294a E. Steinitz. Algebraische Theorie der Körpern, J. de Crelle, v. CXXXVII (1910), pp.~167-309 (new ed.~H. Hasse und R. Baer, Berlin-Leipzig (de Gruyter), 1930).

294b E. Steinitz. Rechteckige Systeme und Moduln in algebraischen Zahlkörpern, Math. Ann., v. LXXI (1912), pp.~328-354 and LXXII (1912), pp.~.297-345.

295 S. Stevin. Les Oeuvres mathématiques . . . , ed.~A. Girard, Leyde (Elzevir), 1634.

296 E. Stiefel. Ueber eine Beziehung zwischen geschlossenen Lie'sche Gruppen und diskontinuierlichen Bewegungsgruppen euklidischer Räume und ihre Anwendung auf die Aufzählung der einfachen Lie'sche Gruppen, Comm. Math. Helv., v. XIV (1941-42), pp.~350-380.

297 T. Stieltjes. Recherches sur les fractions continues, Ann. Fac. Sci. de Toulouse, v. VIII (1894), pp.~J.1-J.122.

298 M. Stifel. Arithmetica integra, Nüremberg, 1544.

299 J. Stirling. Lineae tertii ordinis Newtonianae \ldots, Londini, 1717 (new ed., Paris (Duprat), 1797).

300 A. H. Stone. Paracompactness and product spaces, Bull. Amer. Math. Soc., v. LIV (1948), pp.~977-982.

301a M. H. Stone. The theory of representation for Boolean algebras, Trans. Amer. Math. Soc., v. XL (1936), pp.~37-111.

301b M. H. Stone. Applications of the theory of Boolean rings to general topology, Trans. Amer. Math. Soc., v. XLI (1937), pp.~375-481.

301c M. H. Stone. The generalized Weierstrass approximation theorem, Math. Magazine, v. XXI (1948), pp.~167-183 and 237-254.

302a C. Sturm. Sur les équations différentielles linéaires du second ordre, Journ. de Math., (1), v. I (1836), pp.~106-186.

302b C. Sturm. Sur une classe d'équations à différences partielles, Journ. de Math., (1), v. I (1836), pp.~373-444.

303 L. Sylow. Théorèmes sur les groupes de substitutions, Math. Ann., v. V (1872), pp.~584-594.

304 J. J. Sylvester. Collected Mathematical Papers, 4 vol., Cambridge, 1904-1911.

305 B. Taylor. Methodus Incrementorum directa et inversa, Londini, 1715.

306 P. Tchebychef. Sur deux théorèmes relatifs aux probabilités, Acta Math., v. XIV (1890), pp.~305-315 (= Oeuvres, v. II, pp.~481-491).

307 H. Tietze. Über Funktionen die auf einer abgeschlossenen Menge stetig sind, J. de Crelle, v. CXLV (1915), pp.~9-14.

308a J. Tits. Groupes simple et géométries associées, Proc. Intern. Congress Math., Stockholm, 1962, pp.~197-221.

308b J. Tits. Théorème de Bruhat et sous-groupes paraboliques, C. R. Acad. Sci., v. CCLIV (1962), pp.~2910-2912.

308c J. Tits. Algebraic and abstract simple groups, Ann. of Math., v. LXXX (1964), pp.~313-329.

309a O. Toeplitz. Ueber die Auflösung unendlichvieler linearer Gleichungen mit unendlichvielen Unbekannten, Rend. Circ. Mat. Palermo, v. XXVIII (1909), pp.~88-96.

309b O. Toeplitz. Das Verhältnis von Mathematik und Ideenlehre bei Plato, Quellen und Studien zur Geschichte der Math., Abt. B: Studien, v. I (1929), pp.~3-33.

309c O. Toeplitz. Die mathematische Epinomisstelle, Quellen und Studien zur Geschichte der Math., Abt. B: Studien, v. II (1933), pp.~334-346.

310 E. Torricelli. Opere, 4 vol., ed G. Loria and G. Vassura, Faenza (Montanari), 1919.

311 J. Tropfke. Geschichte der Elementar-Mathematik, 7 vol., 2nd ed., Berlin-Leipzig (de Gruyter), 1921-24.

312 N. Tschebotarów. Grundzüge der Galois'schen Theorie (transl. Schwerdtfeger), Groningen (Noordhoff), 1950.

313 A. Tychonoff. Über die topologische Erweiterung von Räumen, Math. Ann., v. CII (1930), pp.~544-561.

314 P. Urysohn. Ueber die Mächtigkeit der zusammenhängenden Mengen, Math. Ann., v. XCIV (1925), pp.~262-295.

315 A. Vandermonde. Mémoire sur la résolution des équations, Hist. de l'Acad. royale des sciences, year 1771, Paris (1774), pp.~365-416.

316 V. S. Varadarajan. Measures on topological spaces (in Russian), Mat. Sbornik, v. LV (1961), pp.~35-100 (English transl.: Amer. Math. Soc. Transl. (2), vol.~48, pp.~161-228).

317a B. L. van der Waerden. Moderne Algebra, 1st ed., 2 vol., Berlin (Springer), 1930-31.

317b B. L. van der Waerden. Die Klassification der einfachen Lieschen Gruppen, Math. Zeitschr., v. XXXVII (1933), pp.~446-462.

317c B. L. van der Waerden. Zenon und die Grundlagenkrise \ldots, Math. Ann., v. CXVII (1940), pp.~141-161.

317d B. L. van der Waerden. Die Arithmetik der Pythagoreer, I, Math. Ann., v. CXX (1947), pp.~127-153.

318 G. Veronese. Fondamenti di geometria, Padova, 1891.

319 Francisci Vietae. Opera mathematica \ldots, Lugduni Batavorum (Elzevir), 1646.

320a G. Vitali. Una proprieta delle funzioni misurabili, R. Ist. Lombardo Sci. Lett. Rend., (2), v. XXXVIII (1905), pp.~599-603.

320b G. Vitali. Sui gruppi di punti e sulla funzioni di variabili reali, Rend. Acc. Sci. di Torino, v. XLIII (1908), pp.~229-236.

321 H. Vogt. Die Entdeckungsgeschichte des Irrationalen nach Plato und andere Quellen des 4. Jahrhunderts, Bibl. Math., (3), v. X (1909), pp.~97-155.

322a V. Volterra. Leçons sur les fonctions de lignes, Paris (Gauthier-Villars), 1913.

322b V. Volterra. Theory of functionals, London-Glasgow (Blackie and sons), 1930.

323 K. von Fritz. The discovery of incommensurability by Hippasus of Metapontium, Ann. of Math., (2), v. XLVI (1945), pp.~242-264.

324a J. von Neumann. Eine Axiomatisierung der Mengenlehre, J. de Crelle, v. CLIV (1925), pp.~219-240.

324b J. von Neumann. Zur Theorie der Darstellung kontinuierlicher Gruppen, Berliner Sitzungsber., 1927, pp.~76-90.

324c J. von Neumann. Die Axiomatisierung der Mengenlehre, Math. Zeitschr., v. XXVII (1928), pp.~669-752.

324d J. von Neumann. Zur Hilbertschen Beweistheorie, Math. Zeitschr., v. XXVI (1927), pp.~1-46.

324e J. von Neumann. Zum Haarschen Mass in topologischen Gruppen, Compos. Math., v. I (1934), pp.~106-114.

324f J. von Neumann. The uniqueness of Haar's measure, Mat. Sbornik., v. I (= XLIII) (1936), pp.~721-734.

324g J. von Neumann. On rings of operators III, Ann. of Math. (2), v. XLI (1940), pp.~94-161.

325 K. G. V. von Staudt. Beiträge zur Geometrie der Lage, Nürnberg, 1856.

326a J. Wallis. Opera Mathematica, 3 vol., Oxoniae, 1693-95.

326b J. Wallis. Logarithmotechnia Nicolai Mercatoris \ldots, Phil. Trans., v. III (1668), pp.~753-759.

327a H. Weber. Untersuchungen über die allgemeine Grundlagen der Galois'schen Gleichungstheorie, Math. Ann., v. XLIII (1893), pp.~521-544.

327b H. Weber. Leopold Kronecker, Math. Ann., v. XLIII (1893), pp.~1-25.

327c H. Weber. Beweis des Satzes, dass jede eigentlich primitive quadratische Form unendlich viele Primzahlen darzustellen fähig is, Math. Ann., v. XX (1882), pp.~301-329.

327d H. Weber. Theorie der Abels'chen Zahlkörper, IV, Acta Math., v. IX (1886), pp.~105-130.

328a J. Maclagan Wedderburn. A theorem on finite algebras, Trans. Amer. Math. Soc., v. VI (1905), pp.~349-352.

328b J. Maclagan Wedderburn. On hypercomplex numbers, Proc. Lond. Math. Soc. (2), v. VI (1908), pp.~77-118.

328c J. Maclagan Wedderburn. A type of primitive algebra, Trans. Amer. Math. Soc., v. XV (1914), pp.~162-166.

328d J. Maclagan Wedderburn. On division algebras, Trans. Amer. Math. Soc., v. XXII (1921), pp.~129-135.

329a K. Weierstrass. Mathematische Werke, 7 vol., Berlin (Mayer und Müller), 1894-1927.

329b K. Weierstrass. Briefe an P. du Bois-Reymond, Acta Math., v. XXXIX (1923), pp.~199-225.

330a A. Weil. Sur les fonctions elliptiques p-adiques, C. R. Acad. Sci., v. CCIII (1936), p.~22.

330b A. Weil. Sur les espaces uniformes et sur la topologie générale (Actual. Scient. et Ind., no 551), Paris (Hermann), 1937.

330c A. Weil. L'intégration dans les groupes topologiques et ses applications (Actual. Scient. et Ind., no 869), Paris (Hermann), 1940.

330d A. Weil. Foundations of algebraic geometry (Amer. Math. Soc. Coll. Publ., v. 29), new York, 1946.

330e A. Weil. Fibre spaces in Algebraic Geometry, (Notes by A. Wallace), Chicago Univ., 1952.

331a H. Weyl. Lettre à I. Schur, Berliner Sitzungsber., 1924, pp.~338-343.

331b H. Weyl. Theorie der Darstellung kontinuierlicher halbeinfacher Gruppen durch lineare Transformationen, Math. Zeitschr., v. XXIII (1925), pp.~271-309, v. XXIV (1926), pp.~328-395 and 789-791.

331c H. Weyl. Symmetry, Princeton (Princeton Univ. Press), 1952.

332 H. Weyl und F. Peter. Die Vollständigkeit der primitiven Darstellungen einer geschlossenen kontinuierlichen Gruppe, Math. Ann., v. XCVII (1927), pp.~737-755.

333 A. N. Whitehead. On cardinal numbers, Amer. Journ. of Math., v. XXIV (1902), pp.~367-394.

334a J. H. C. Whitehead. On the decomposition of an infinitesimal group, Proc. Camb. Phil. Soc., v. XXXII (1936), pp.~229-237.

334b J. H. C. Whitehead. Certain equations in the algebra of a semi-simple infinitesimal group, Quart. Journ. of Math., (2), v. VIII (1937), pp.~220-237.

335 H. Wieleitner. Der ``Tractatus de latitudinibus formarum'' des Oresme, Bibl. Math. (3), v. XIII (1912), pp.~115-145.

336 N. Wiener. Differential space, J. Math. Phys. M. I. T., v. II (1923), pp.~131-174.

337a E. Witt. Theorie der quadratischen Formen in beliebigen Körpern, J. de Crelle, v. CLXXVI (1937), pp.~31-44.

337b E. Witt. Treue Darstellung Lieschen Ringe, J. de Crelle, v. CLXXVII (1937), pp.~152-160.

337c E. Witt. Spiegelungsgruppen und Aufzählung halbeinfacher Liescher Ringe, Abh. math. Sem. Univ. Hamburg, v. XIV (1941), pp.~289-322.

338 E. Wright. Table of Latitudes, 1599.

339a W. H. Young. On upper and lower integration, Proc. London Math. Soc. (2), v. II (1905), pp.~52-66.

339b W. H. Young. A new method in the theory of integration, Proc. London Math. Soc. (2), v. IX (1911), pp.~15-50.

340a O. Zariski. The compactness of the Riemann manifold of an abstract field of algebraic functions, Bull. Amer. Math. Soc., v. L (1944), pp.~683-691.

340b O. Zariski. Generalized local rings, Summa Bras. Math., v. I (1946), pp.~169-195.

341 H. Zassenhaus. Lehrbuch der Gruppentheorie, Bd. I, Leipzig-Berlin (Teubner), 1937.

342a E. Zermelo. Beweis dass jede Menge wohlgeordnet werden kann, Math. Ann., v. LIX (1904), pp.~514-516.

342b E. Zermelo. Neuer Beweis für die Möglichkeit einer Wohlordnung, Math. Ann., v. LXV (1908), pp.~107-128.

342c E. Zermelo. Untersuchung über die Grundlagen der Mengenlehre, Math. Ann., v. LXV (1908), pp.~261-281.

343 G. Zolotareff. Sur la théorie des nombres complexes, Journ. de Math. (3), v. VI (1880), pp.~51-84 and 129-166.

344 M. Zorn. A remark on method in transfinite algebra, Bull. Amer. Math. Soc., v. XLI (1935), pp.~667-670.

345 A. Zygmund. Trigonometrical series, 2 vol., Cambridge (Univ. Press), 1959.
